\chapter{Clocks}
\label{ch:clocks}

SIRCUS has one oscillator. Every other rate in the system is that oscillator divided down, so every timing in the
console is a whole number of cycles of it, and the three chips stay in step without any of them watching another.

\section{The Master Clock}
\label{sec:master-clock}

The master clock is the system oscillator. Its frequency depends on the region, because the video signal that SIRCADA
makes must match the colour subcarrier of the television standard:

\begin{table}[H]
	\centering
	\begin{tabular}{lll}
		\toprule
		\textbf{Region} & \textbf{Master clock} & \textbf{Derivation}                 \\
		\midrule
		NTSC            & 21,477,272 Hz         & The NTSC colour subcarrier times 6  \\
		PAL             & 21,281,370 Hz         & The PAL colour subcarrier times 4.8 \\
		\bottomrule
	\end{tabular}
	\caption{The Master Clock}
	\label{tab:master-clock}
\end{table}

The oscillator itself, its tolerance and its drive circuit are not yet specified.

\section{The Clock Rail}
\label{sec:clock-rail}

The system divides the master clock by two, and the result is \texttt{CLKI}. All three chips take the same
\texttt{CLKI}: the CPU, SIRCADA and SIRCAS each have a single clock input and no other clock. \texttt{CLKI} is
10,738,636 Hz in an NTSC system and 10,640,685 Hz in a PAL system. The CPU does one execution phase on each rising edge
(\hwmanual{}, signal descriptions chapter), SIRCADA does one PPU cycle, and SIRCAS does one cycle of its own. The three
devices therefore count the same cycles at the same moments, and a cycle count in one manual is a cycle count in the
others.

\begin{table}[H]
	\centering
	\begin{tabular}{llp{6.6cm}}
		\toprule
		\textbf{Chip} & \textbf{What one \texttt{CLKI} cycle is} & \textbf{Longer intervals}                                     \\
		\midrule
		SIRC-1        & One execution phase                      & An instruction is 6 cycles                                    \\
		SIRCADA       & One PPU cycle                            & A pixel is 2 cycles; a line is 682 cycles                     \\
		SIRCAS        & One cycle of the device                  & A sample is 336 cycles; a bit clock period is 10 or 11 cycles \\
		\bottomrule
	\end{tabular}
	\caption{What the Clock Means to Each Chip}
	\label{tab:clock-meaning}
\end{table}

The phase relationship between the rising edges at the three chips, and the electrical characteristics of the clock,
are not yet specified. The figures in this manual show logical timing only, as the manuals of the chips do.

\section{Rates in Both Regions}
\label{sec:rates}

Every figure below follows from the clock and the cycle counts of Table~\ref{tab:clock-meaning}.

\begin{table}[H]
	\centering
	\begin{tabular}{lrr}
		\toprule
		                                               & \textbf{NTSC} & \textbf{PAL}  \\
		\midrule
		Master clock                                   & 21,477,272 Hz & 21,281,370 Hz \\
		\texttt{CLKI}                                  & 10,738,636 Hz & 10,640,685 Hz \\
		\texttt{CLKI} period                           & 93.1 ns       & 94.0 ns       \\
		CPU instruction (6 cycles)                     & 558.7 ns      & 563.9 ns      \\
		Pixel (2 cycles)                               & 186.2 ns      & 188.0 ns      \\
		Line (682 cycles)                              & 63.51 $\mu$s  & 64.09 $\mu$s  \\
		Lines in a frame                               & 262           & 312           \\
		Frame (cycles)                                 & 178,684       & 212,784       \\
		Frame                                          & 16.64 ms      & 20.00 ms      \\
		Frame rate                                     & 60.10 Hz      & 50.01 Hz      \\
		Vertical blanking (cycles)                     & 25,916        & 60,016        \\
		Vertical blanking                              & 2.41 ms       & 5.64 ms       \\
		Audio sample (336 cycles)                      & 31.29 $\mu$s  & 31.58 $\mu$s  \\
		Audio sample rate                              & 31,960.2 Hz   & 31,668.7 Hz   \\
		Envelope update (21,504 cycles)                & 2.002 ms      & 2.021 ms      \\
		Converter bit clock (mean, \texttt{CLKI}/10.5) & 1,022,727 Hz  & 1,013,399 Hz  \\
		\bottomrule
	\end{tabular}
	\caption{Rates in Both Regions}
	\label{tab:system-rates}
\end{table}

\section{NTSC and PAL}
\label{sec:regions}

A PAL console is the same console with a different oscillator. Nothing else in the logic changes, and a program written
for NTSC runs on PAL. Three things differ:

\begin{itemize}
	\item \textbf{The picture.} The frame has 312 lines instead of 262, so vertical blanking is longer, 5.64~ms
	      against 2.41~ms, and the frame rate is 50 Hz against 60 Hz. The region of the display adapter is fixed when
	      it is made and the program reads it from the \texttt{OUT} bit of the Status register
	      (\videomanual{}, registers chapter).
	\item \textbf{The CPU.} It runs 0.9\% slower, and a frame holds more of its instructions, about 35,464 against
	      29,781, since the frame is longer. Code that is timed in instructions to fill a frame needs a figure for each
	      region.
	\item \textbf{The sound.} SIRCAS has no region input. It runs at the \texttt{CLKI} it is given, so every note in a
	      PAL console is 15.9 cents flat of an NTSC console, and a sequence takes 0.9\% longer. A program that must be in
	      tune in both regions builds its pitch tables for the clock of the region it finds itself in
	      (\audiomanual{}, timing chapter).
\end{itemize}

\section{The Range SIRCAS Accepts}

SIRCAS works at any steady \texttt{CLKI} from 9.0 to 12.0 MHz (\audiomanual{}, timing chapter). Both regions are well
inside the range, and the clock circuit of any other system built around the chip may divide its own oscillator to any
frequency in it. The CPU and SIRCADA have no stated range; they are specified at the \texttt{CLKI} of this system.
